\documentclass[12pt]{article}
\usepackage[a4paper,margin=1in]{geometry}
\usepackage[T1]{fontenc}
\usepackage{lmodern,microtype}
\usepackage{amsmath,amssymb,amsthm,mathtools}
\usepackage{booktabs,array,enumitem,xcolor}
\usepackage[numbers,sort&compress]{natbib}
\usepackage[colorlinks=true,linkcolor=blue!55!black,citecolor=blue!55!black,urlcolor=blue!55!black]{hyperref}
\linespread{1.07}

\newtheorem{theorem}{Theorem}[section]
\newtheorem{proposition}[theorem]{Proposition}
\newtheorem{corollary}[theorem]{Corollary}
\theoremstyle{remark}
\newtheorem{remark}[theorem]{Remark}
\newcommand{\norm}[1]{\left\lVert#1\right\rVert}

\title{A Complement-Resolvent Budget Audit\\
\large The Norm-Only Oblique Route Fails, Leaving Validated Contour Inverses as the Next Wall}
\author{Liang Wang}
\date{July 2026}

\begin{document}
\maketitle

\begin{abstract}
RH-189 gives an exact oblique Feshbach determinant factorization, but the
physical complement resolvent remains unknown.  This paper tests the
strongest elementary norm-only substitute.

For a biorthogonal packet with oblique projector $P=VW^*$ and condition
$\chi=\norm P$, let $Q=I-P$.  In orthonormal complement coordinates, the
Feshbach block obeys the sharp elementary bound
\[
 \norm D\le\norm Q\norm A=\chi\norm A.
\]
If a cycle root has radius $\rho$ and a contour has radius $\delta$, a
Neumann resolvent estimate would require
\[
 \rho-\delta>\chi\norm A.
\]
This criterion is intentionally optimistic: it ignores all nonnormal
cancellation and uses the best ideal packet contour factor.

Applied to all 126 RH-185 windows with contour radius $0.4$ times the root
half-spacing, the norm-only clearance is negative everywhere.  The complement
operator bound ranges from $71.46$ to $6.40\times10^5$ while the
minimum contour modulus is at most $0.719$.  No norm-only complement
resolvent exists, no full norm-only certificate passes, and even an
unrealistically unit complement-resolvent Schur test has zero successes.

This is a useful negative result: it rules out a cheap norm closure, not the
physical Feshbach route.  The next required object is a validated contour
inverse exploiting the actual complement spectrum, as in the finite-mesh and
operator-ball machinery of RH-167--168.  No physical Riesz shell, Gate A,
Hilbert--Polya, or RH result is claimed.
\end{abstract}

\section{The complement problem after Feshbach factorization}

RH-189 writes the physical operator in oblique packet/complement coordinates
as
\begin{equation}\label{eq:block}
 S^{-1}AS=\begin{pmatrix}K&B\\C&D\end{pmatrix}
\end{equation}
and gives
\begin{equation}\label{eq:feshbach}
 \det(zI-A)=\det(zI-D)
 \det\left(zI-K-B(zI-D)^{-1}C\right).
\end{equation}
The coupling product from RH-188 controls $B(zI-D)^{-1}C$ only after a bound
on $(zI-D)^{-1}$ is available \citep{WangRH188,WangRH189}.

One might first try the universal estimate
\begin{equation}\label{eq:universal}
 \norm{(zI-D)^{-1}}
 \le\frac1{|z|-\norm D}.
\end{equation}
This paper shows exactly how far that shortcut is from the physical data.

\section{Oblique complement norm bound}

Let $P=VW^*$ with $W^*V=I$ and $Q=I-P$.  By RH-186,
$\norm Q=\norm P=\chi$ for a nontrivial oblique decomposition.  The
orthonormal complement construction of RH-189 chooses an isometry $Z$ onto
$\operatorname{Ran}Q=\ker W^*$ and gives
$D=Z^*QAZ$ \citep{WangRH189}.

\begin{theorem}[Universal oblique complement bound]\label{thm:complement}
For every bounded $A$,
\begin{equation}\label{eq:D-bound}
 \norm D\le\norm Q\norm A\le\chi\norm A.
\end{equation}
If $\rho>\delta>0$ and
\begin{equation}\label{eq:clearance}
 \rho-\delta>\chi\norm A,
\end{equation}
then every $z$ with $|z-\zeta|=\delta$ and $|\zeta|=\rho$ satisfies
\begin{equation}\label{eq:resolvent}
 \norm{(zI-D)^{-1}}
 \le\frac1{\rho-\delta-\chi\norm A}.
\end{equation}
\end{theorem}

\begin{proof}
The representation $D=Z^*QAZ$ and the isometry of $Z$ give
$\norm D\le\norm Q\norm A$.  On the contour,
$|z|\ge|\zeta|-|z-\zeta|=\rho-\delta$.  If \eqref{eq:clearance} holds,
\[
 zI-D=z\left[I-z^{-1}D\right]
\]
has a convergent Neumann inverse, and summing its geometric series gives
\eqref{eq:resolvent}.
\end{proof}

The theorem is a sufficient condition only.  It does not claim that the
complement has spectrum reaching $\norm D$; it records what a norm-only
proof would need to establish.

\section{Contour choice}

For a length-$L$ cycle of radius $\rho$, the root half-spacing is
\begin{equation}\label{eq:spacing}
 s_L=\rho\sin(\pi/L).
\end{equation}
We choose the deliberately conservative contour radius
\begin{equation}\label{eq:contour}
 \delta=0.4s_L<s_L.
\end{equation}
This leaves root shells geometrically disjoint and is more generous than a
small perturbative contour.  The packet resolvent factor is given its best
ideal estimate $a=1/\delta$ for the optimistic Schur comparison.

For each RH-185 window, the stored source-cycle radius supplies $\rho$, the
stored length supplies $L$, and the stored operator norm and oblique condition
supply $\norm A$ and $\chi$.

\section{Audit results}

There are 126 windows.  The norm-only complement data are:
\begin{center}
\begin{tabular}{lrrr}
\toprule
quantity & minimum & median & maximum\\
\midrule
$\chi\norm A$ & $71.46$ & $1.5547\times10^3$
& $6.3996\times10^5$\\
$\rho-\delta$ & $0.3913$ & $0.5260$ & $0.7188$\\
clearance & $-6.3996\times10^5$ & $-1.5542\times10^3$
& $-70.979$\\
\bottomrule
\end{tabular}
\end{center}
Thus
\begin{equation}\label{eq:no-clearance}
 \#\{\text{norm-only complement resolvent successes}\}=0.
\end{equation}
The result holds in particular for all 38 local $\sigma=0.01,L=4$ windows.

The optimistic comparison sets the complement resolvent factor to one,
which is strictly better than the unavailable norm-only bound.  Even then,
the product
\begin{equation}\label{eq:optimistic}
 a\,b\,c
\end{equation}
has minimum $1.1116$ and zero successes.  Inserting the actual norm-only
resolvent would only worsen it.

\section{Interpretation}

The negative result has three layers:
\begin{enumerate}[leftmargin=2em]
\item the orthonormal-coordinate complement bound still carries one full
 oblique factor $\chi$;
\item the physical $\chi$ values are too large for radial norm separation;
\item even an artificially favorable complement factor does not close the
 smallest observed coupling product on the chosen contour.
\end{enumerate}

None of these layers proves that the actual complement spectrum is large.
Nonnormal matrices can have resolvents much smaller than a norm-only
Neumann bound at selected points, or much larger near pseudospectral regions.
The correct next calculation must use actual sample inverses and outward
operator balls.

\section{Validated contour route}

The RH-167--168 architecture provides the appropriate replacement:
\begin{enumerate}[leftmargin=2em]
\item choose a finite contour mesh around each candidate root;
\item compute a nominal inverse at each mesh point;
\item bound inverse defects and operator displacement;
\item use the Banach denominator to cover the continuous contour;
\item feed the resulting $d(z)$ into the directed product
 $a(z)d(z)bc$.
\end{enumerate}
This route may exploit structure invisible to $\chi\norm A$.  It is more
expensive and is now the precise physical $D$ leaf rather than a vague
``resolvent estimate'' request.

The finite validation step can be stated as a direct Banach-lemma theorem.

\begin{proposition}[Mesh-and-operator-ball inverse certificate]\label{prop:validated-inverse}
Let $D_0$ be a nominal complement block, let
$\norm{D-D_0}\le\varepsilon$, and choose a mesh point $z_0$ with a computed
matrix $M_0$.  Suppose
\begin{equation}\label{eq:inverse-defect}
 \norm{I-M_0(z_0I-D_0)}\le\eta.
\end{equation}
For every $z$ with $|z-z_0|\le h$, define
\begin{equation}\label{eq:q-bound}
 q=\eta+\norm{M_0}(h+\varepsilon).
\end{equation}
If $q<1$, then $zI-D$ is invertible and
\begin{equation}\label{eq:validated-resolvent}
 \norm{(zI-D)^{-1}}\le\frac{\norm{M_0}}{1-q}.
\end{equation}
\end{proposition}

\begin{proof}
Expanding around the nominal mesh equation gives
\[
 I-M_0(zI-D)
 =I-M_0(z_0I-D_0)-M_0\bigl((z-z_0)I-(D-D_0)\bigr).
\]
Its norm is at most $q$.  If $q<1$, the Banach lemma makes
$M_0(zI-D)$ invertible and bounds its inverse by $(1-q)^{-1}$.  Multiplying
by $M_0$ yields \eqref{eq:validated-resolvent}.
\end{proof}

This proposition specifies the data missing from the present archive:
nominal complement matrices, outward operator radii, contour mesh radii,
inverse residuals, and inverse norms.  It also explains why the failure of a
radial norm disk is not decisive.  A good local inverse $M_0$ may exist far
inside the disk $|z|\le\norm D$.

\section{Quantitative distance to the norm-only route}

The best elementary clearance is still about $-70.98$, whereas the relevant
contour moduli are below one.  Thus replacing the ambient $\chi^2$ estimate
by the sharp orthonormal-coordinate factor $\chi$ improves the bound by
orders of magnitude but leaves a large finite gap.  This corrected audit is
important: the negative conclusion is not an artifact of paying an
unnecessary second oblique factor.

The optimistic unit-complement comparison is logically separate.  Its
minimum $1.1116$ says that on the chosen contour even a complement inverse
of norm one would not close the current product.  A successful exact audit
must therefore gain jointly from the actual packet resolvent geometry,
directional coupling coordinates, and complement inverse; improving only
the radial complement norm estimate is insufficient.

\section{Data contract for the next experiment}

The next archive should store the nominal matrices $K,B,C,D$ for every
surviving window, not merely their norms.  It should also store contour
centers, mesh points, mesh covering radii, approximate inverses, inverse
defects, and outward operator radii.  From these primitive quantities one
can independently reconstruct \eqref{eq:q-bound}, the continuous
complement bound, and the final Schur product.

Predeclaring this data contract prevents a posteriori contour movement.  A
contour that meets the complement spectrum is a finite rejection of that
root shell; it should not be silently shrunk until a favorable inverse is
found.  Conversely, a positive margin must survive the full mesh cells and
operator ball, not just the sampled contour points.

\section{Boundary}

This paper proves the norm-only complement criterion and audits its complete
failure on the current 126 finite windows.  It does not establish a
complement spectral obstruction, exclude validated contour inverses, prove a
Schur/Riesz certificate, or close any all-level interface.  The local
biorthogonal candidate remains an exploratory seed only.

\bibliographystyle{plainnat}
\bibliography{references}
\end{document}
